% GENERATED FILE. Edit canonical lessons, then run npm run generate:books.
% canonical-source-hash: fnv1a64:d1a97ef310154099
% canonical-lessons: HI-C445-kisti, HI-C445-mallah, HI-C445-jal, HI-C445-jvar, HI-C445-sip

\chapter{Boats, Nets and Shells}
\label{ch:hi-boats-nets-and-shells}
\hlchaptermodality{445}

\begin{chapteropening}
\textbf{By the end of this chapter:} \emph{I can talk about a small boat, the boatman and his net, the tide and shells on the shore.}

Complete the last lesson of chapter 445: I can talk about a small boat, the boatman and his net, the tide and shells on the shore.
\end{chapteropening}

\section[\texorpdfstring{\hi{किश्ती} (kiśtī)}{किश्ती (kiśtī)}]{\hi{किश्ती} (kiśtī) — a small boat, a skiff}

\label{lesson:HI-C445-kisti}



\begin{quote}
Before the new one: say the Hindi for a coast, then the Hindi for a dip.
\end{quote}

\subsection*{You'll want to know: \hi{किश्ती}}
\textbf{\hi{किश्ती}} — \emph{kiśtī} — \textquotedblleft{}a small boat, a skiff\textquotedblright{}.

\begin{quote}
\hi{झील} \hi{में} \hi{एक} \hi{छोटी} \hi{किश्ती} \hi{तैर} \hi{रही} \hi{थी।}
\end{quote}

\begin{quote}
\emph{jhīl meṁ ek choṭī kiśtī tair rahī thī.}
\end{quote}

\begin{quote}
A little boat was floating on the lake.
\end{quote}

Feminine.

\begin{cousinweb}
\textbf{\hi{किश्ती}} (\emph{kiśtī}) is usually traced to Persian \emph{kištī}, \textquotedblleft{}a boat, a ship\textquotedblright{}; songs and poems love it.
\end{cousinweb}

\begin{grammarlens}[title={What you've built}]
One more word for getting around.
\end{grammarlens}

\subsection*{Your turn}
\begin{itemize}
\raggedright
  \item \emph{Say it:} \emph{kiśtī}
  \item \emph{Say it:} \emph{kiśtī}, once more
  \item \emph{Say it:} \emph{jhīl meṁ ek choṭī kiśtī tair rahī thī.} — A little boat was floating on the lake.
  \item \emph{Recall:} say the Hindi for a coast, then the Hindi for a dip, then say \emph{kiśtī} again
\end{itemize}

\begin{tcolorbox}[breakable,colback=teal!4,colframe=teal!35!black,title={Before you move on}]
What does \hi{किश्ती} mean? (\textquotedblleft{}A small boat, a skiff\textquotedblright{}.) Say it once more.
\end{tcolorbox}
\section[\texorpdfstring{\hi{मल्लाह} (mallāh)}{मल्लाह (mallāh)}]{\hi{मल्लाह} (mallāh) — a boatman}

\label{lesson:HI-C445-mallah}



\begin{quote}
Before the new one: say the Hindi for a dip, then the Hindi for a small boat.
\end{quote}

\subsection*{You'll want to know: \hi{मल्लाह}}
\textbf{\hi{मल्लाह}} — \emph{mallāh} — \textquotedblleft{}a boatman\textquotedblright{}.

\begin{quote}
\hi{मल्लाह} \hi{ने} \hi{हमें} \hi{नदी} \hi{के} \hi{पार} \hi{पहुँचाया।}
\end{quote}

\begin{quote}
\emph{mallāh ne hamẽ nadī ke pār pahũcāyā.}
\end{quote}

\begin{quote}
The boatman took us across the river.
\end{quote}

Masculine.

\begin{cousinweb}
\textbf{\hi{मल्लाह}} (\emph{mallāh}) is usually traced to Arabic \emph{mallāḥ}, \textquotedblleft{}a sailor\textquotedblright{}, through Persian; on Indian rivers it is the man who rows the ferry.
\end{cousinweb}

\begin{grammarlens}[title={What you've built}]
One more word for getting around.
\end{grammarlens}

\subsection*{Your turn}
\begin{itemize}
\raggedright
  \item \emph{Say it:} \emph{mallāh}
  \item \emph{Say it:} \emph{mallāh}, once more
  \item \emph{Say it:} \emph{mallāh ne hamẽ nadī ke pār pahũcāyā.} — The boatman took us across the river.
  \item \emph{Recall:} say the Hindi for a dip, then the Hindi for a small boat, then say \emph{mallāh} again
\end{itemize}

\begin{tcolorbox}[breakable,colback=teal!4,colframe=teal!35!black,title={Before you move on}]
What does \hi{मल्लाह} mean? (\textquotedblleft{}A boatman\textquotedblright{}.) Say it once more.
\end{tcolorbox}
\section[\texorpdfstring{\hi{जाल} (jāl)}{जाल (jāl)}]{\hi{जाल} (jāl) — a net}

\label{lesson:HI-C445-jal}



\begin{quote}
Before the new one: say the Hindi for a small boat, then the Hindi for a boatman.
\end{quote}

\subsection*{You'll want to know: \hi{जाल}}
\textbf{\hi{जाल}} — \emph{jāl} — \textquotedblleft{}a net\textquotedblright{}.

\begin{quote}
\hi{मल्लाह} \hi{ने} \hi{पानी} \hi{में} \hi{जाल} \hi{फेंका।}
\end{quote}

\begin{quote}
\emph{mallāh ne pānī meṁ jāl phẽkā.}
\end{quote}

\begin{quote}
The boatman cast his net into the water.
\end{quote}

Masculine.

\begin{cousinweb}
\textbf{\hi{जाल}} (\emph{jāl}) is a tadbhava of Sanskrit \emph{jāla}, \textquotedblleft{}a net\textquotedblright{} (CDIAL 5213), which lists Hindi \emph{jāl}, a net for birds or fish; it is a different word from \emph{jala}, water.
\end{cousinweb}

\begin{grammarlens}[title={What you've built}]
One more word for getting around.
\end{grammarlens}

\subsection*{Your turn}
\begin{itemize}
\raggedright
  \item \emph{Say it:} \emph{jāl}
  \item \emph{Say it:} \emph{jāl}, once more
  \item \emph{Say it:} \emph{mallāh ne pānī meṁ jāl phẽkā.} — The boatman cast his net into the water.
  \item \emph{Recall:} say the Hindi for a small boat, then the Hindi for a boatman, then say \emph{jāl} again
\end{itemize}

\begin{tcolorbox}[breakable,colback=teal!4,colframe=teal!35!black,title={Before you move on}]
What does \hi{जाल} mean? (\textquotedblleft{}A net\textquotedblright{}.) Say it once more.
\end{tcolorbox}
\section[\texorpdfstring{\hi{ज्वार} (jvār)}{ज्वार (jvār)}]{\hi{ज्वार} (jvār) — the high tide, a flood tide}

\label{lesson:HI-C445-jvar}



\begin{quote}
Before the new one: say the Hindi for a boatman, then the Hindi for a net.
\end{quote}

\subsection*{You'll want to know: \hi{ज्वार}}
\textbf{\hi{ज्वार}} — \emph{jvār} — \textquotedblleft{}the high tide, a flood tide\textquotedblright{}.

\begin{quote}
\hi{शाम} \hi{को} \hi{ज्वार} \hi{आता} \hi{है} \hi{और} \hi{पानी} \hi{ऊपर} \hi{चढ़ता} \hi{है।}
\end{quote}

\begin{quote}
\emph{śām ko jvār ātā hai aur pānī ūpar caṛhtā hai.}
\end{quote}

\begin{quote}
In the evening the tide comes in and the water rises.
\end{quote}

Masculine.

\begin{cousinweb}
\textbf{\hi{ज्वार}} (\emph{jvār}) goes back to \emph{*javacāra}, \textquotedblleft{}a flood tide\textquotedblright{}, a form CDIAL reconstructs (CDIAL 5167), which lists Hindi \emph{jvār}. A different word with the same spelling names the millet sorghum.
\end{cousinweb}

\begin{grammarlens}[title={What you've built}]
One more word for getting around.
\end{grammarlens}

\subsection*{Your turn}
\begin{itemize}
\raggedright
  \item \emph{Say it:} \emph{jvār}
  \item \emph{Say it:} \emph{jvār}, once more
  \item \emph{Say it:} \emph{śām ko jvār ātā hai aur pānī ūpar caṛhtā hai.} — In the evening the tide comes in and the water rises.
  \item \emph{Recall:} say the Hindi for a boatman, then the Hindi for a net, then say \emph{jvār} again
\end{itemize}

\begin{tcolorbox}[breakable,colback=teal!4,colframe=teal!35!black,title={Before you move on}]
What does \hi{ज्वार} mean? (\textquotedblleft{}The high tide, a flood tide\textquotedblright{}.) Say it once more.
\end{tcolorbox}
\section[\texorpdfstring{\hi{सीप} (sīp)}{सीप (sīp)}]{\hi{सीप} (sīp) — a shell (of an oyster or a mussel)}

\label{lesson:HI-C445-sip}



\begin{quote}
Before the new one: say the Hindi for a net, then the Hindi for the high tide.
\end{quote}

\subsection*{You'll want to know: \hi{सीप}}
\textbf{\hi{सीप}} — \emph{sīp} — \textquotedblleft{}a shell (of an oyster or a mussel)\textquotedblright{}.

\begin{quote}
\hi{मीरा} \hi{को} \hi{किनारे} \hi{पर} \hi{एक} \hi{सुंदर} \hi{सीप} \hi{मिली।}
\end{quote}

\begin{quote}
\emph{Mīrā ko kināre par ek sundar sīp milī.}
\end{quote}

\begin{quote}
Mira found a pretty shell on the shore.
\end{quote}

Feminine.

\begin{cousinweb}
\textbf{\hi{सीप}} (\emph{sīp}) goes back to \emph{*sippī}, \textquotedblleft{}a bivalve shell\textquotedblright{}, a form CDIAL reconstructs (CDIAL 13417), which lists Hindi \emph{sīp}: the shell in which a pearl grows.
\end{cousinweb}

\begin{grammarlens}[title={What you've built}]
One more word for getting around.
\end{grammarlens}

\subsection*{Your turn}
\begin{itemize}
\raggedright
  \item \emph{Say it:} \emph{sīp}
  \item \emph{Say it:} \emph{sīp}, once more
  \item \emph{Say it:} \emph{Mīrā ko kināre par ek sundar sīp milī.} — Mira found a pretty shell on the shore.
  \item \emph{Recall:} say the Hindi for a net, then the Hindi for the high tide, then say \emph{sīp} again
\end{itemize}

\begin{tcolorbox}[breakable,colback=teal!4,colframe=teal!35!black,title={Before you move on}]
What does \hi{सीप} mean? (\textquotedblleft{}A shell (of an oyster or a mussel)\textquotedblright{}.) Say it once more.
\end{tcolorbox}
